\section{Pruning inside an alternating ancestor}\label{sec:alternating}

We use the same failed-chart construction to force the next almost simple group inside a high-degree alternating ancestor. The successive anchors become transitive, primitive, non-affine and non-Cartesian as the exclusions proceed.

\begin{convention}[Prepared alternating configuration]\label{c:alternating-configuration}
Use the pattern, prior-anchor and full-stabilizer-label conventions of the structure exclusions. Thus the present almost simple ancestor quotient is \(A_v\le M\le S_v\) of high degree on \(\Lambda\); labels under test and witnesses are in the interval above \(H\). For a witness \(K_i>H\), \([H,K_i]\) has the prepared link properties, including (\(*\)): a normal subgroup is retained in \(H\) or supplemented by \(H\) (missing), soluble normals are retained, and two commuting normals cannot both be missing (in the core quotient they would both contain the nonabelian monolith). Where earlier exclusions are used we have the previously good \(H<E<K_i\) common within a pattern, with link properties at \(E\) as well. Patterns can have a whole line of \(q\) distinct labels with \(K_i\cap K_j\le K_l\) for distinct \(i,j\) and all \(l\) on the line; cross-chart pairs when invoked have \(K_i\cap K_j=H\).
\end{convention}

\subsection{Permutation orbit bounds and primitivity}

\begin{lemma}[Two alternating permutation images]\label{c:alternating-orbits}
If a group acts on two finite sets, inducing at least the alternating group on each, then its number of orbits on ordered pairs of pairs, one coordinate from each set in each pair, is bounded absolutely.
\end{lemma}

\begin{proof}
We use the following orbit bound for two finite set actions on sizes \(s,t\), each inducing at least the alternating group on that set. The number of orbits on ordered pairs of pairs (one coordinate from each set per pair) is absolutely bounded. Indeed the joint action is subdirect on the two induced groups. In large degrees either the product of the alternating groups is contained, or the two alternating simples are linked diagonally, hence the same degree and natural identification up to relabeling. This follows by the normal kernels of subdirect projections and simplicity (normal subgroups not containing the alternating simple must centralize it); use also the usual automorphism theorem for alternating groups in high degrees. If only one degree is large the large alternating factor acts independently by the same argument. Equality patterns and the bounded small sizes give the assertion.
\end{proof}

\begin{proposition}[Subsets and equal-block partitions]\label{c:alternating-primitivity}
In the prepared alternating configuration, full stabilizers of proper nonempty subsets and nontrivial equal-block-size partitions are excluded. Subsequent good anchors are primitive.
\end{proposition}

\begin{proof}
First exclude stabilizers of proper nonempty subsets. For a cross-chart
pair, the intersection contains supported alternating groups on at most
four cells. Thus $H$ has boundedly many point orbits, and every subset
whose stabilizer is a witness is a union of these orbits. Only boundedly
many subsets, and hence witness labels, are possible.

Next exclude nontrivial equal-block-size partitions, which suffice for
primitivity once transitivity is known. Write $s$ for the number of
blocks and $k$ for their common size. In a witness label, the
intersection $Y$ of $M$ with the base block group is normal. If $Y$ is
retained and $k\ge3$, every individually supported $A_k$ lies in $H$.
Its primitivity makes the restriction of any other tested
$H$-invariant equivalence to that block either universal or discrete.
In the discrete case a point has an equivalent point outside the block,
since the other partition also has block size at least two. The
supported transitive group fixes that outside point and would make it
equivalent to every point of the original block, a contradiction.
Consequently each original block lies in a block of the other partition.
Applying this in both directions to two retained systems with block
sizes at least three forces their equality. For retained $k=2$, the
even simultaneous flips on any two blocks lie in $H$. A distinct
matching fails invariance by flipping a block while leaving one of its
points' outside partner fixed, using the high degree.

If $HY=K_i$, then $H$ induces the full symmetric top action on the
blocks. For a cross-chart pair of missing partitions,
\Cref{c:alternating-orbits} bounds the orbits on ordered pairs of block
cells. Moreover $H$ contains the supported alternating groups on the
points of each cell, so each such orbit has boundedly many point-pair
orbits above it, including when both points lie in the same cell.
Thus $H$ has boundedly many orbits on $\Lambda^2$. Every invariant
equivalence is a union of those orbits, giving the absolute partition
bound used in \eqref{c:partition-bound}. This is too small for all the
distinct pattern labels. Here and below the chart field size exceeds
these absolute bounds, after coloring by the indicated finitely many
cases.
\end{proof}

\begin{proposition}[Elementary abelian regular groups]\label{c:alternating-affine}
For a primitive previously good anchor, normalizers of elementary abelian regular groups are excluded.
\end{proposition}

\begin{proof}
Exclude next normalizers of elementary abelian regular groups. Such groups lie in \(A_v\) (parity of translations, high degree) and would thus be retained soluble normals. They would be minimal normal in the primitive previously good \(E\), by transitivity of every nontrivial normal subgroup and regularity. Distinct ones must commute, impossible for regular self-centralizing groups.
\end{proof}

\subsection{Cartesian decompositions}

\begin{lemma}[Retained Cartesian products]\label{c:cartesian-retained}
For a primitive anchor avoiding the affine exclusion, a normalized nontrivial Cartesian decomposition has transitive slots of common size \(k\ge5\). Its full stabilizer has a normal product \(Y=A_k^s\le A_v\) and induces at least \(A_s\) on slots. Retained products determine at most one system in a pattern.
\end{lemma}

\begin{proof}
Next exclude stabilizers of nontrivial Cartesian decompositions of the point set (product coordinates \(\Lambda\simeq\prod_{a=1}^s \Lambda_a\), \(s\ge2,\ |\Lambda_a|\ge2\); stabilize their system, i.e. equality-of-coordinate partitions, allowing slot permutations). Primitivity of \(E\) gives slot transitivity (an orbit of slots otherwise gives an invariant nontrivial equivalence), and thus common slot set size \(k\), say on \(s\ge2\) slots. Sizes \(2,3,4\) would give canonical elementary abelian regular groups (\(S_2,A_3,V_4\), respectively, multiplied over slots). For \(k\ge5\) let \(Y_i=A_k^s\) acting slotwise, inside \(A_v\) and normal in the label \(K_i\). The induced slot action of \(K_i\) contains \(A_s\) by commutators of true slot permutations. These stabilizers are proper (Cartesian coordinate agreement prevents full 2-transitivity). If \(Y_i\le H\), it is minimal normal in \(E\), and has trivial centralizer on \(\Lambda\) (point stabilizer in \(Y_i\) self-normalizing). Thus all retained ones in the pattern must coincide, determining the slots themselves by the orbits of factor subproducts.
\end{proof}

\begin{lemma}[Missing normal powers with primitive factor action]\label{c:primitive-product-interval}
Let \([H,K_i]\) be a prepared interval and \(Y_i\trianglelefteq K_i\) a direct power of nonabelian simple groups with \(HY_i=K_i\). If the action of \(H\) on its simple factors is primitive, then \(H\cap Y_i\) is a nontrivial coordinate product, and every group between \(H\) and \(K_i\) intersects \(Y_i\) in a coordinate product.
\end{lemma}

\begin{proof}
Here and in the next exclusion we use a simple consequence of \Cref{f:prop-interval-cases}. If \(HY_i=K_i\) for a direct normal power of nonabelian simples and \(H\)'s action on the factors is primitive, then \(P=H\cap Y_i\) is a nontrivial coordinate product and all groups between \(H,K_i\) intersect \(Y_i\) in coordinate products. Indeed \(Y_i\) is minimal normal, missed by the base so embeds as the minimal normal modulo core, and the exact interval is that of invariant subgroups above \(P\). In the earlier reductions, the subdirect case and the all-nonbottom-subdirect case are now both impossible: ties must be factor-index invariant, so primitivity allows no chain of comparable distinct proper nontrivial subdirect groups. This leaves precisely the asserted product case.
\end{proof}

\begin{lemma}[Common Cartesian refinement]\label{c:cartesian-refinement}
A finite pattern line of missing Cartesian labels has a common refinement into finite coordinate sets. It is constructed from the tuple intersections of the slotwise products \(E\cap Y_i\). The original Cartesian systems are recovered by grouping the fine coordinates, and \(E\) permutes the fine coordinate directions.
\end{lemma}

\begin{proof}
For missing Cartesian labels \(HY_i=K_i\), \(H\) induces the full top image, transitive and containing \(A_s\), hence primitive on slot indices. On a line of these labels write
\[
 N_i=E\cap Y_i=\prod_a R_{i,a}\trianglelefteq E
\]
a slotwise product strictly above $H\cap Y_i$. In
$\bar E=E/\operatorname{core}_E H$, each $\bar N_i$ contains the
nonabelian simple-power monolith $W$. A simple factor $C$ of $W$ lies
in one of the $\bar R_{i,a_i}$ for each $i$: the commutators of those
slot groups with $W$ generate $W$, are normal in $W$, and commute for
distinct slots.

Apply the two commutator steps from the proof of
\Cref{c:tensor-refinement}. First, since the $N_i$ are normal in $E$
and their images contain the perfect group $W$, their intersection
$N=\bigcap_iN_i$ still has image containing $W$. Next,
$[N,R_{i,a_i}]\le N\cap R_{i,a_i}$ has image containing $C$.
The subgroups $N\cap R_{i,a_i}$ are normal in $N$, so iterating
commutators among them shows that the actual tuple intersection
\[
 T_\gamma=\bigcap_iR_{i,a_i},\qquad \gamma=(a_i),
\]
has image containing $C$.

Take these intersections for all slot tuples. They are permuted by
$E$ and generate a nontrivial normal subgroup, which is transitive
because $E$ is primitive. Distinct tuple groups commute, since they
act in distinct slots of at least one original system. Consider their
action through the abstract product $\prod_\gamma T_\gamma$ and fix
a point. If a product of tuple elements fixes that point, retain only
the components assigned to one prescribed slot $a$ of an original
system $i$. This partial product still fixes the point: it acts only
on coordinate $a$, where its action agrees with that of the original
fixing product. Repeating this restriction through every original
system isolates one tuple component at a time. Thus the stabilizer
in the abstract product is exactly the product of the individual
point stabilizers.

Transitivity now identifies the point set with the product of the
individual coset sets. Each original coordinate depends only on the
tuple components assigned to its slot. The maps from these separate
batches to the original coordinate sets have a bijective product map,
so each is itself a bijection. Grouping the fine coordinates therefore
recovers the original equality-of-coordinate partitions. Discard the
trivial coset sets; the remaining coordinate directions are still
permuted by $E$.
\end{proof}

\begin{proposition}[Completion of Cartesian pruning]\label{c:cartesian-pruning}
Missing Cartesian labels are excluded. Together with \Cref{c:cartesian-retained}, this excludes all nontrivial Cartesian decompositions.
\end{proposition}

\begin{proof}
Fix two distinct systems $i,j$ and group the fine coordinates by the
occupied pairs of their original slot indices. On an aggregate set of
size $u$, the supported $A_u$ lies in $K_i\cap K_j$, hence in every
label on the pattern line. If another system split that aggregate,
then on a slice with all outside coordinates fixed its nonconstant
coordinate partitions would give a nontrivial Cartesian decomposition
invariant under the natural $A_u$. Such a split has $u\ge4$ and
distinguishes pairs agreeing in some coordinate from pairs agreeing
in none, contrary to 2-transitivity. Thus there is only one nontrivial
fine coordinate over each occupied pair.

The map from fine coordinates to occupied pairs is $H$-equivariant:
$H\le E$ permutes the tuple groups and their coordinate directions,
and stabilizes both original systems. By
\Cref{c:alternating-orbits}, the two slot actions bound the number of
ordered-pair orbits on this set. Each original system is recovered
by an $H$-invariant grouping of the fine coordinates, as proved in
\Cref{c:cartesian-refinement}. Distinct full stabilizer labels therefore
give distinct invariant partitions. The bound
\eqref{c:partition-bound}, with fibers of size one, contradicts the
number of labels on the pattern line.
\end{proof}

\subsection{Diagonal actions and the almost simple conclusion}

\begin{lemma}[Primitive groups after Cartesian pruning]\label{c:primitive-diagonal}
Let \(E_0\) be a primitive group satisfying the preceding affine and Cartesian exclusions. If it is not almost simple, it normalizes a transitive group \(Y=B^s\), with \(B\) nonabelian simple and \(s\ge2\), whose point stabilizer is a full possibly twisted diagonal. Its degree is \(|B|^{s-1}\). This includes the pair of opposed regular simple groups.
\end{lemma}

\begin{proof}
Recall a short primitive-group observation after these exclusions. A minimal normal \(P=C^u\) of a primitive group \(E_0\) here is transitive with \(C\) nonabelian simple (an abelian one would be regular), \(E_0=P(E_0)_x\) at a point \(x\). Its point stabilizer \(Q=P_x\) has no proper intermediate increase towards \(P\) normalized by \((E_0)_x\). Taking its coordinate product closure, \(Q\) is either a product or subdirect. In the product case \(u>1\) would give a Cartesian coset product; in the subdirect case the same holds for multiple tie parts (uniform sizes by factor transitivity of \((E_0)_x\)). Indeed the point stabilizer of \(E_0\) is transitive on factors (minimal normality, with \(P\) itself not permuting them). In the product situation, including when taking batches given by multiple tie parts, the cosets \(P/Q\) accordingly split as a product of coset sets. Their sizes are uniform and nontrivial, \(P\) respects this product coordinatewise, and \((E_0)_x\) respects it by its conjugation/transport of the factors or batches with their respective stabilizers in \(Q\). Thus for \(u>1\) we have just a full diagonal, possibly with twists, and \(v=|C|^{u-1}\). For \(u=1\) with nontrivial \(C_{E_0}(P)\), that centralizer must be transitive so both are regular opposite simple groups, again giving a transitive power \(C^2\) with full diagonal point stabilizer. Consequently any non-almost-simple case here normalizes a transitive \(Y=B^s,\ s\ge2\), \(B\) nonabelian simple, with full diagonal point stabilizer.
\end{proof}

\begin{proposition}[Diagonal normalizers]\label{c:diagonal-pruning}
Under the successively primitive, non-affine and non-Cartesian anchor hypotheses, the full normalizers of the groups \(Y\) in \Cref{c:primitive-diagonal} are excluded. Both retained and missing normal powers are covered.
\end{proposition}

\begin{proof}
Finally exclude the normalizers \(K_i\) of all such \(Y\). They are proper since \(Y\le A_v\) by perfection, and induce at least \(A_s\) on factors by diagonal-coordinate factor permutations and commutators. For retained \(Y\le H\), transitivity of \(E\) on factors, when present, makes \(Y\) minimal \(E\)-normal with trivial centralizer (the diagonal is self-normalizing), leaving no distinct such system. If instead \(E\) has multiple factor orbits, their products are each transitive so use at least \(s-1\) factors each, forcing just two opposed regular simple groups, the only minimal normals in \(E\), again determining \(Y\). If \(K_i\)'s top action is intransitive, likewise \(s=2\), and at least one individual factor is retained by (\(*\)); such factors are regular simple minimal normals of \(E\), determining \(Y\) by taking also their opposite regular group (distinct such minimal normals commute so determine the same pair).

\textit{The missing transitive case.} It remains that \(HY=K_i\) and \(H\) induces a transitive top action containing \(A_s\). By the interval product consequence, \(N=E\cap Y=\prod R_a\) is a nontrivial slot product with each projection proper (since \(E<K_i\)). Choose a minimal \(E\)-normal \(P'=C^u\) inside it. Each factor of \(P'\) is contained actually in one \(R_a\), since \([P',R_a]\le P'\cap R_a\), normal in \(P'\), and those commutators generate \(P'\) by perfection. This assignment is \(E\)-equivariant with some common positive number \(h\) per slot, so \(u=sh\). The primitive observation applied to \(E\) gives \(Q=P'_x\) a full diagonal. If \(h>1\), take inside \(P'\) the product of the projections of \(Q\) along the factor blocks of the assignment. It is a proper intermediate increase, normalized by \(E_x\), impossible. If \(h=1\), point degrees give \(|C|^{s-1}=|B|^{s-1}\), contrary to containment in the proper \(R_a\).
\end{proof}

\begin{theorem}[The next alternating-ancestor socle]\label{c:alternating-conclusion}
After the exclusions, every tested descendant with its prepared interval is primitive almost simple. Its actual simple socle is transitive, lies in \(A_v\), and is supplemented by \(H\).
\end{theorem}

\begin{proof}
These arguments use only successively primitive, non-affine, and non-Cartesian prior anchors as needed and finite case subdivisions. Thus after the exclusions the tested descendant with its prepared interval is primitive almost simple, with an actual simple transitive socle in \(A_v\), supplemented by \(H\) via (\(*\)) (it cannot be retained, since then factoring by the base core leaves a soluble quotient by Schreier, inconsistent with the prepared interval), as desired.
\end{proof}
